\color{unidarkgreen}{\section[(обязательное) Уставки функций РЗА]{(обязательное)\\Уставки функций РЗА}\label{app:set}}
\color{black}



%%%%%%%%%%%%%%%%%%%%%%%%%%%%%%%%%%%%%%%%%%%%%%%%%%%%%%%%%% МТЗ %%%%%%%%%%%%%%%%%%%%%%%%%%%%%%%%%%%%%%%%%%%%%%%%%%%%%%%%%%%%%%%%%%
\par
В таблице \ref{mtzvn:tbl1} приведены параметры, необходимые для настройки МТЗ.

\footnotesize
\begin{longtable}{|>{\centering\arraybackslash}m{5.3cm}|>{\centering\arraybackslash}m{3.3cm}|>{\centering\arraybackslash}m{4.2cm}|>{\centering\arraybackslash}m{1.8cm}|>{\centering\arraybackslash}m{1cm}|}
\caption{Параметры для настройки МТЗ\hfill\vspace{-0.5\baselineskip}}\label{mtzvn:tbl1}\\ 
\hline
\rowcolor{gray!20}
Параметр (Параметр на ИЧМ) & Условное обозначение на схеме & Значение/ Диапазон & Единица измерения & Шаг \\ 
\hline
\endfirsthead
%\caption{\emph{Продолжение\hfill\vspace{-0.5\baselineskip}}} \\ % Отступ обозначения таблицы от таблицы и выравнивание надписи слева
\caption*{\hspace{3pt}\emph{Продолжение таблицы \ref{mtzvn:tbl1}\hfill\vspace{-0.5\baselineskip}}} \\ % сделано по ГОСТ 2.105 п.6.8.7
\hline
\rowcolor{gray!20}
Параметр (Параметр на ИЧМ) & Условное обозначение на схеме & Значение/ Диапазон & Единица измерения & Шаг \\ 
%\hline
\endhead
%\hline
\endfoot
%\hline
\endlastfoot
%==+t1*LLN0|LVTTOC> Общие уставки
\multicolumn{5}{|c|}{ МТЗ 1 ст } \\ \hline 
\centering Ввод функции в работу (Ввод\_функции)  & \centering SGF1 & \centering Не предусмотрено \\ Предусмотрено & \centering -- & \centering \arraybackslash -- \\
\hline
\centering Ток срабатывания (Iср)  & \centering I> & \centering 0,10...30,00 & \centering о.е. & \centering \arraybackslash 0,01 \\
\hline
\centering Ток срабатывания грубого органа в режиме управляющего напряжения (Iср\_груб)  & \centering -- & \centering 0,10...30,00 & \centering о.е. & \centering \arraybackslash 0,01 \\
\hline
\centering Тип пуска по напряжению (Тип\_КПОН)  & \centering SGF2 & \centering Управляющее напряжение \\ Вольтметровая блокировка & \centering -- & \centering \arraybackslash -- \\
\hline
\centering Режим контроля от БНТ (Реж\_БНТ)  & \centering SGF3 & \centering Не предусмотрено \\ Предусмотрено & \centering -- & \centering \arraybackslash -- \\
\hline
\centering Режим КПОН при неисправности ЦН (Реж\_БНН\_КПОН)  & \centering SGF4 & \centering Деблокировка \\ Блокировка & \centering -- & \centering \arraybackslash -- \\
\hline
\centering Режим контроля от КПОН (Реж\_КПОН)  & \centering SGF5 & \centering Не предусмотрено \\ Предусмотрено & \centering -- & \centering \arraybackslash -- \\
\hline
\centering Режим контроля СВ НН (Контр\_СВ)  & \centering SGF6 & \centering Не предусмотрено \\ Блокировка ступени при включенном СВ \\ Блокировка ступени при отключенном СВ & \centering -- & \centering \arraybackslash -- \\
\hline
\centering Независимая выдержка времени срабатывания (Tср)  & \centering T1 & \centering 0...100000 & \centering мс & \centering \arraybackslash 10 \\
\hline
\multicolumn{5}{|c|}{ МТЗ 2 ст } \\ \hline 
\centering Ввод функции в работу (Ввод\_функции)  & \centering SGF1 & \centering Не предусмотрено \\ Предусмотрено & \centering -- & \centering \arraybackslash -- \\
\hline
\centering Ток срабатывания (Iср)  & \centering I> & \centering 0,10...30,00 & \centering о.е. & \centering \arraybackslash 0,01 \\
\hline
\centering Ток срабатывания грубого органа в режиме управляющего напряжения (Iср\_груб)  & \centering -- & \centering 0,10...30,00 & \centering о.е. & \centering \arraybackslash 0,01 \\
\hline
\centering Тип пуска по напряжению (Тип\_КПОН)  & \centering SGF2 & \centering Управляющее напряжение \\ Вольтметровая блокировка & \centering -- & \centering \arraybackslash -- \\
\hline
\centering Режим контроля от БНТ (Реж\_БНТ)  & \centering SGF3 & \centering Не предусмотрено \\ Предусмотрено & \centering -- & \centering \arraybackslash -- \\
\hline
\centering Режим КПОН при неисправности ЦН (Реж\_БНН\_КПОН)  & \centering SGF4 & \centering Деблокировка \\ Блокировка & \centering -- & \centering \arraybackslash -- \\
\hline
\centering Режим контроля от КПОН (Реж\_КПОН)  & \centering SGF5 & \centering Не предусмотрено \\ Предусмотрено & \centering -- & \centering \arraybackslash -- \\
\hline
\centering Режим контроля СВ НН (Контр\_СВ)  & \centering SGF6 & \centering Не предусмотрено \\ Блокировка ступени при включенном СВ \\ Блокировка ступени при отключенном СВ & \centering -- & \centering \arraybackslash -- \\
\hline
\centering Независимая выдержка времени срабатывания (Tср)  & \centering T1 & \centering 0...100000 & \centering мс & \centering \arraybackslash 10 \\
\hline
\multicolumn{5}{|c|}{ МТЗ 3 ст } \\ \hline 
\centering Ввод функции в работу (Ввод\_функции)  & \centering SGF1 & \centering Не предусмотрено \\ Предусмотрено & \centering -- & \centering \arraybackslash -- \\
\hline
\centering Ток срабатывания (Iср)  & \centering I> & \centering 0,10...30,00 & \centering о.е. & \centering \arraybackslash 0,01 \\
\hline
\centering Ток срабатывания грубого органа в режиме управляющего напряжения (Iср\_груб)  & \centering -- & \centering 0,10...30,00 & \centering о.е. & \centering \arraybackslash 0,01 \\
\hline
\centering Тип пуска по напряжению (Тип\_КПОН)  & \centering SGF2 & \centering Управляющее напряжение \\ Вольтметровая блокировка & \centering -- & \centering \arraybackslash -- \\
\hline
\centering Режим контроля от БНТ (Реж\_БНТ)  & \centering SGF3 & \centering Не предусмотрено \\ Предусмотрено & \centering -- & \centering \arraybackslash -- \\
\hline
\centering Режим КПОН при неисправности ЦН (Реж\_БНН\_КПОН)  & \centering SGF4 & \centering Деблокировка \\ Блокировка & \centering -- & \centering \arraybackslash -- \\
\hline
\centering Режим контроля от КПОН (Реж\_КПОН)  & \centering SGF5 & \centering Не предусмотрено \\ Предусмотрено & \centering -- & \centering \arraybackslash -- \\
\hline
\centering Режим контроля СВ НН (Контр\_СВ)  & \centering SGF6 & \centering Не предусмотрено \\ Блокировка ступени при включенном СВ \\ Блокировка ступени при отключенном СВ & \centering -- & \centering \arraybackslash -- \\
\hline
\centering Независимая выдержка времени срабатывания (Tср)  & \centering T1 & \centering 0...100000 & \centering мс & \centering \arraybackslash 10 \\
\hline
\multicolumn{5}{|c|}{ КПОН } \\ \hline 
\centering Напряжение срабатывания (Uср)  & \centering U< & \centering 2,0...100,0 & \centering В & \centering \arraybackslash 0,1 \\
\hline
\centering Напряжение срабатывания ОП (U2ср)  & \centering U2> & \centering 2,0...50,0 & \centering В & \centering \arraybackslash 0,1 \\
\hline
\centering Режим пуска (Реж\_пуска)  & \centering SGF1 & \centering По Uмин \\ Комбинированный \\ Внешний & \centering -- & \centering \arraybackslash -- \\
\hline
\multicolumn{5}{|c|}{ БНТ } \\ \hline 
\centering Отношение тока второй гармоники к току основной гармоники (Ih2/Ih1)  & \centering I2г/I1г> & \centering 5...100 & \centering \% & \centering \arraybackslash 1 \\
\hline
\centering Ток блокировки (Iср)  & \centering I> & \centering 0,1...15,0 & \centering о.е. & \centering \arraybackslash 0,1 \\
\hline
\centering Перекрестная блокировка (Перекрест\_блок)  & \centering SGF1 & \centering Не предусмотрено \\ Предусмотрено & \centering -- & \centering \arraybackslash -- \\
\hline
\multicolumn{5}{|c|}{ БЛЗШ } \\ \hline 
\centering Выбор ступени блокировки (БлокЛЗШ\_выбор\_ст)  & \centering SGF1 & \centering Не предусмотрено \\ 1 ступень \\ 2 ступень \\ 3 ступень & \centering -- & \centering \arraybackslash -- \\
\hline
\multicolumn{5}{|c|}{ МТЗ } \\ \hline 
\centering Сборка токовых цепей (Сборка\_ток\_цепей)  & \centering SGF1 & \centering Звезда \\ Треугольник & \centering -- & \centering \arraybackslash -- \\
\hline
%===t1
\end{longtable}
\normalsize


%%%%%%%%%%%%%%%%%%%%%%%%%%%%%%%%%%%%%%%%%%%%%%%%%%%%%%%%%% ТО %%%%%%%%%%%%%%%%%%%%%%%%%%%%%%%%%%%%%%%%%%%%%%%%%%%%%%%%%%%%%%%%%%
\par

В таблице \ref{to:tbl1} приведены параметры, необходимые для настройки ТО.

\footnotesize
\begin{longtable}{|>{\centering\arraybackslash}m{5.3cm}|>{\centering\arraybackslash}m{3.3cm}|>{\centering\arraybackslash}m{4.2cm}|>{\centering\arraybackslash}m{1.8cm}|>{\centering\arraybackslash}m{1cm}|}
\caption{Параметры для настройки ТО\hfill\vspace{-0.5\baselineskip}}\label{to:tbl1}\\ 
\hline
\rowcolor{gray!20}
Параметр (Параметр на ИЧМ) & Условное обозначение на схеме & Значение/ Диапазон & Единица измерения & Шаг \\ 
\hline
\endfirsthead
\caption*{\hspace{3pt}\emph{Продолжение таблицы \ref{to:tbl1}\hfill\vspace{-0.5\baselineskip}}} \\ % сделано по ГОСТ 2.105 п.6.8.7
\hline
\rowcolor{gray!20}
Параметр (Параметр на ИЧМ) & Условное обозначение на схеме & Значение/ Диапазон & Единица измерения & Шаг \\ 
%\hline
\endhead
%\hline
\endfoot
%\hline
\endlastfoot
%==+t1*PTOC1|T_LVTOC
\centering Ввод функции в работу (Ввод\_функции)  & \centering SGF1 & \centering Не предусмотрено \\ Предусмотрено & \centering -- & \centering \arraybackslash -- \\
\hline
\centering Сборка токовых цепей (Сборка\_ток\_цепей)  & \centering SGF2 & \centering Звезда \\ Треугольник & \centering -- & \centering \arraybackslash -- \\
\hline
\centering Ток срабатывания (Iср)  & \centering I> & \centering 0,10...30,00 & \centering о.е. & \centering \arraybackslash 0,01 \\
\hline
\centering Выдержка времени срабатывания (Tср)  & \centering T1 & \centering 0...10000 & \centering мс & \centering \arraybackslash 10 \\
\hline
%===t1
\end{longtable}
\normalsize


%%%%%%%%%%%%%%%%%%%%%%%%%%%%%%%%%%%%%%%%%%%%%%%%%%%%%%%%%% ЛЗТ %%%%%%%%%%%%%%%%%%%%%%%%%%%%%%%%%%%%%%%%%%%%%%%%%%%%%%%%%%%%%%%%%%
\par

В таблице \ref{lzt:tbl1} приведены параметры, необходимые для настройки ЛЗТ.

\footnotesize
\begin{longtable}{|>{\centering\arraybackslash}m{5.3cm}|>{\centering\arraybackslash}m{3.3cm}|>{\centering\arraybackslash}m{4.2cm}|>{\centering\arraybackslash}m{1.8cm}|>{\centering\arraybackslash}m{1cm}|}
\caption{Параметры для настройки ЛЗТ\hfill\vspace{-0.5\baselineskip}}\label{lzt:tbl1}\\ 
\hline
\rowcolor{gray!20}
Параметр (Параметр на ИЧМ) & Условное обозначение на схеме & Значение/ Диапазон & Единица измерения & Шаг \\ 
\hline
\endfirsthead
\caption*{\hspace{3pt}\emph{Продолжение таблицы \ref{lzt:tbl1}\hfill\vspace{-0.5\baselineskip}}} \\ % сделано по ГОСТ 2.105 п.6.8.7
\hline
\rowcolor{gray!20}
Параметр (Параметр на ИЧМ) & Условное обозначение на схеме & Значение/ Диапазон & Единица измерения & Шаг \\ 
%\hline
\endhead
%\hline
\endfoot
%\hline
\endlastfoot
%==+t1*PTRC1|TTOCLGC_UIRZ
\centering Ввод функции в работу (Ввод\_функции)  & \centering SGF1 & \centering Не предусмотрено \\ Предусмотрено & \centering -- & \centering \arraybackslash -- \\
\hline
\centering Пуск от МТЗ 2 ступени (Пуск\_МТЗ\_ст2)  & \centering SGF2 & \centering Не предусмотрено \\ Предусмотрено & \centering -- & \centering \arraybackslash -- \\
\hline
\centering Пуск от МТЗ 3 ступени (Пуск\_МТЗ\_ст3)  & \centering SGF3 & \centering Не предусмотрено \\ Предусмотрено & \centering -- & \centering \arraybackslash -- \\
\hline
\centering Выдержка времени срабатывания (Tср)  & \centering T1 & \centering 0,00...2,00 & \centering с & \centering \arraybackslash 0,01 \\
\hline
%===t1
\end{longtable}
\normalsize


%%%%%%%%%%%%%%%%%%%%%%%%%%%%%%%%%%%%%%%%%%%%%%%%%%%%%%%%%% ЗП %%%%%%%%%%%%%%%%%%%%%%%%%%%%%%%%%%%%%%%%%%%%%%%%%%%%%%%%%%%%%%%%%%
\par
В таблице \ref{zp:tbl1} приведены параметры, необходимые для настройки ЗП.

\footnotesize
\begin{longtable}{|>{\centering\arraybackslash}m{5.3cm}|>{\centering\arraybackslash}m{3.3cm}|>{\centering\arraybackslash}m{4.2cm}|>{\centering\arraybackslash}m{1.8cm}|>{\centering\arraybackslash}m{1cm}|}
\caption{Параметры для настройки ЗП\hfill\vspace{-0.5\baselineskip}}\label{zp:tbl1}\\ 
\hline
\rowcolor{gray!20}
Параметр (Параметр на ИЧМ) & Условное обозначение на схеме & Значение/ Диапазон & Единица измерения & Шаг \\ 
\hline
\endfirsthead
\caption*{\hspace{3pt}\emph{Продолжение таблицы \ref{zp:tbl1}\hfill\vspace{-0.5\baselineskip}}} \\ % сделано по ГОСТ 2.105 п.6.8.7
\hline
\rowcolor{gray!20}
Параметр (Параметр на ИЧМ) & Условное обозначение на схеме & Значение/ Диапазон & Единица измерения & Шаг \\ 
%\hline
\endhead
%\hline
\endfoot
%\hline
\endlastfoot
%==+t1*HVPTOC1|T_AUTOVC
\centering Ввод функции в работу (Ввод\_функции)  & \centering SGF1 & \centering Не предусмотрено \\ Предусмотрено & \centering -- & \centering \arraybackslash -- \\
\hline
\centering Ток срабатывания (Iср)  & \centering I> & \centering 0,10...5,00 & \centering о.е. & \centering \arraybackslash 0,01 \\
\hline
\centering Выдержка времени срабатывания (Tср)  & \centering T1 & \centering 0...3600000 & \centering мс & \centering \arraybackslash 100 \\
\hline
%===t1
\end{longtable}
\normalsize

%%%%%%%%%%%%%%%%%%%%%%%%%%%%%%%%%%%%%%%%%%%%%%%%%%%%%%%%%% ЗОП %%%%%%%%%%%%%%%%%%%%%%%%%%%%%%%%%%%%%%%%%%%%%%%%%%%%%%%%%%%%%%%%%%
\par
В таблице \ref{zop:tbl1} приведены параметры, необходимые для настройки ЗОП.

\footnotesize
\begin{longtable}{|>{\centering\arraybackslash}m{5.3cm}|>{\centering\arraybackslash}m{3.3cm}|>{\centering\arraybackslash}m{4.2cm}|>{\centering\arraybackslash}m{1.8cm}|>{\centering\arraybackslash}m{1cm}|}
\caption{Параметры для настройки ЗОП\hfill\vspace{-0.5\baselineskip}}\label{zop:tbl1}\\ 
\hline
\rowcolor{gray!20}
Параметр (Параметр на ИЧМ) & Условное обозначение на схеме & Значение/ Диапазон & Единица измерения & Шаг \\ 
\hline
\endfirsthead
\caption*{\hspace{3pt}\emph{Продолжение таблицы \ref{zop:tbl1}\hfill\vspace{-0.5\baselineskip}}} \\ % сделано по ГОСТ 2.105 п.6.8.7
\hline
\rowcolor{gray!20}
Параметр (Параметр на ИЧМ) & Условное обозначение на схеме & Значение/ Диапазон & Единица измерения & Шаг \\ 
%\hline
\endhead
%\hline
\endfoot
%\hline
\endlastfoot
%==+t1*NSPTOC1|LVNSTOC
\centering Ввод функции в работу (Ввод\_функции)  & \centering SGF1 & \centering Не предусмотрено \\ Предусмотрено & \centering -- & \centering \arraybackslash -- \\
\hline
\centering Ток срабатывания ОП (I2ср)  & \centering I2> & \centering 0,04...4,00 & \centering о.е. & \centering \arraybackslash 0,01 \\
\hline
\centering Отношение тока ОП к току ПП (Kнесим)  & \centering I2/I1> & \centering 0,02...1,00 & \centering о.е. & \centering \arraybackslash 0,01 \\
\hline
\centering Режим работы ЗОП (Реж\_ЗОП)  & \centering SGF2 & \centering по I2 \\ по I2/I1 \\ по I2 или по I2/I1 & \centering -- & \centering \arraybackslash -- \\
\hline
\centering Выдержка времени срабатывания (Tср)  & \centering T1 & \centering 100...20000 & \centering мс & \centering \arraybackslash 10 \\
\hline
%===t1
\end{longtable}
\normalsize


%%%%%%%%%%%%%%%%%%%%%%%%%%%%%%%%%%%%%%%%%%%%%%%%%%%%%%%%%% ТК ЗДЗ %%%%%%%%%%%%%%%%%%%%%%%%%%%%%%%%%%%%%%%%%%%%%%%%%%%%%%%%%%%%%%%%%%
\par

В таблице \ref{tkzdz:tbl1} приведены параметры, необходимые для настройки ТК ЗДЗ.

\footnotesize
\begin{longtable}{|>{\centering\arraybackslash}m{5.3cm}|>{\centering\arraybackslash}m{3.3cm}|>{\centering\arraybackslash}m{4.2cm}|>{\centering\arraybackslash}m{1.8cm}|>{\centering\arraybackslash}m{1cm}|}
\caption{Параметры для настройки ТК ЗДЗ\hfill\vspace{-0.5\baselineskip}}\label{tkzdz:tbl1}\\ 
\hline
\rowcolor{gray!20}
Параметр (Параметр на ИЧМ) & Условное обозначение на схеме & Значение/ Диапазон & Единица измерения & Шаг \\ 
\hline
\endfirsthead
\caption*{\hspace{3pt}\emph{Продолжение таблицы \ref{tkzdz:tbl1}\hfill\vspace{-0.5\baselineskip}}} \\ % сделано по ГОСТ 2.105 п.6.8.7
\hline
\rowcolor{gray!20}
Параметр (Параметр на ИЧМ) & Условное обозначение на схеме & Значение/ Диапазон & Единица измерения & Шаг \\ 
%\hline
\endhead
%\hline
\endfoot
%\hline
\endlastfoot
%==+t1*PTOC1|T_LVARCTOC
\centering Ввод функции в работу (Ввод\_функции)  & \centering SGF1 & \centering Не предусмотрено \\ Предусмотрено & \centering -- & \centering \arraybackslash -- \\
\hline
\centering Ток срабатывания (Iср)  & \centering I> & \centering 0,10...30,00 & \centering о.е. & \centering \arraybackslash 0,01 \\
\hline
\centering Выбор пускового органа (Выбор\_пуск)  & \centering SGF2 & \centering Внутренний \\ МТЗ 1 ст \\ МТЗ 2 ст \\ МТЗ 3 ст & \centering -- & \centering \arraybackslash -- \\
\hline
%===t1
\end{longtable}
\normalsize


%%%%%%%%%%%%%%%%%%%%%%%%%%%%%%%%%%%%%%%%%%%%%%%%%%%%%%%%%% ТО РПН %%%%%%%%%%%%%%%%%%%%%%%%%%%%%%%%%%%%%%%%%%%%%%%%%%%%%%%%%%%%%%%%%%

\par
В таблице \ref{torpn:tbl1} приведены параметры, необходимые для настройки ТО РПН.

\footnotesize
\begin{longtable}{|>{\centering\arraybackslash}m{5.3cm}|>{\centering\arraybackslash}m{3.3cm}|>{\centering\arraybackslash}m{4.2cm}|>{\centering\arraybackslash}m{1.8cm}|>{\centering\arraybackslash}m{1cm}|}
\caption{Параметры для настройки ТО РПН\hfill\vspace{-0.5\baselineskip}}\label{torpn:tbl1}\\ 
\hline
\rowcolor{gray!20}
Параметр (Параметр на ИЧМ) & Условное обозначение на схеме & Значение/ Диапазон & Единица измерения & Шаг \\ 
\hline
\endfirsthead
\caption*{\hspace{3pt}\emph{Продолжение таблицы \ref{torpn:tbl1}\hfill\vspace{-0.5\baselineskip}}} \\ % сделано по ГОСТ 2.105 п.6.8.7
\hline
\rowcolor{gray!20}
Параметр (Параметр на ИЧМ) & Условное обозначение на схеме & Значение/ Диапазон & Единица измерения & Шаг \\ 
%\hline
\endhead
%\hline
\endfoot
%\hline
\endlastfoot
%==+t1*PTOC1|LTCBLKTOC
\centering Ввод функции в работу (Ввод\_функции)  & \centering SGF1 & \centering Не предусмотрено \\ Предусмотрено & \centering -- & \centering \arraybackslash -- \\
\hline
\centering Ток срабатывания (Iср)  & \centering I> & \centering 0,10...30,00 & \centering о.е. & \centering \arraybackslash 0,01 \\
\hline
%===t1
\end{longtable}
\normalsize

%%%%%%%%%%%%%%%%%%%%%%%%%%%%%%%%%%%%%%%%%%%%%%%%%%%%%%%%%% КЦН НН %%%%%%%%%%%%%%%%%%%%%%%%%%%%%%%%%%%%%%%%%%%%%%%%%%%%%%%%%%%%%%%%%%

\par

В таблице \ref{kznnn:tbl1} приведены параметры, необходимые для настройки КЦН НН.

\small
\begin{longtable}{|>{\centering\arraybackslash}m{5.3cm}|>{\centering\arraybackslash}m{3.3cm}|>{\centering\arraybackslash}m{4.2cm}|>{\centering\arraybackslash}m{1.8cm}|>{\centering\arraybackslash}m{1cm}|}
\caption{Параметры для настройки КЦН НН\hfill\vspace{-0.5\baselineskip}}\label{kznnn:tbl1}\\ 
\hline
\rowcolor{gray!20}
Параметр (Параметр на ИЧМ) & Условное обозначение на схеме & Значение/ Диапазон & Единица измерения & Шаг \\ 
\hline
\endfirsthead
%\caption{\emph{Продолжение\hfill\vspace{-0.5\baselineskip}}} \\ % Отступ обозначения таблицы от таблицы и выравнивание надписи слева
\caption*{\hspace{3pt}\emph{Продолжение таблицы \ref{kznnn:tbl1}\hfill\vspace{-0.5\baselineskip}}} \\ % сделано по ГОСТ 2.105 п.6.8.7
\hline
\rowcolor{gray!20}
Параметр (Параметр на ИЧМ) & Условное обозначение на схеме & Значение/ Диапазон & Единица измерения & Шаг \\ 
%\hline
\endhead
%\hline
\endfoot
%\hline
\endlastfoot
%==+t1*RVTR1|LVRBVTR
\multicolumn{5}{|c|}{ КЦН } \\ \hline 
\centering Ввод функции в работу (Ввод\_функции)  & \centering SGF1 & \centering Не предусмотрено \\ Предусмотрено & \centering -- & \centering \arraybackslash -- \\
\hline
\centering Режим пуска (Реж\_пуска)  & \centering SGF2 & \centering Внутренний \\ Внешний & \centering -- & \centering \arraybackslash -- \\
\hline
\centering Напряжение срабатывания (Uср)  & \centering U< & \centering 2...100 & \centering В & \centering \arraybackslash 1 \\
\hline
\centering Напряжение срабатывания ОП (U2ср)  & \centering U2> & \centering 2,0...50,0 & \centering В & \centering \arraybackslash 0,1 \\
\hline
\centering Выдержка времени срабатывания (Tср)  & \centering T1 & \centering 0...100000 & \centering мс & \centering \arraybackslash 10 \\
\hline
%===t1
\end{longtable}
\normalsize



%%%%%%%%%%%%%%%%%%%%%%%%%%%%%%%%%%%%%%%%%%%%%%%%%%%%%%%%%% ЛО ГЗоткл %%%%%%%%%%%%%%%%%%%%%%%%%%%%%%%%%%%%%%%%%%%%%%%%%%%%%%%%%%%%%%%%%%
\par

В таблице \ref{logzotkl:tbl1} приведены параметры, необходимые для настройки ЛО ГЗоткл.

\footnotesize
\begin{longtable}{|>{\centering\arraybackslash}m{5.3cm}|>{\centering\arraybackslash}m{3.3cm}|>{\centering\arraybackslash}m{4.2cm}|>{\centering\arraybackslash}m{1.8cm}|>{\centering\arraybackslash}m{1cm}|}
\caption{Параметры для настройки логики отключения при срабатывании отключающего контакта газового реле трансформатора\hfill\vspace{-0.5\baselineskip}}\label{logzotkl:tbl1}\\ 
\hline
\rowcolor{gray!20}
Параметр (Параметр на ИЧМ) & Условное обозначение на схеме & Значение/ Диапазон & Единица измерения & Шаг \\ 
\hline
\endfirsthead
\caption*{\hspace{3pt}\emph{Продолжение таблицы \ref{logzotkl:tbl1}\hfill\vspace{-0.5\baselineskip}}} \\ % сделано по ГОСТ 2.105 п.6.8.7
\hline
\rowcolor{gray!20}
Параметр (Параметр на ИЧМ) & Условное обозначение на схеме & Значение/ Диапазон & Единица измерения & Шаг \\ 
%\hline
\endhead
%\hline
\endfoot
%\hline
\endlastfoot
%==+t1*PTRC1|TTRGASLGC
\multicolumn{5}{|c|}{ БУ } \\ \hline 
\centering Выдержка времени срабатывания на блокировку (Тбл)  & \centering T1 & \centering 0...30000 & \centering мс & \centering \arraybackslash 10 \\
\hline
\multicolumn{5}{|c|}{ ЛО } \\ \hline 
\centering Ввод функции в работу (Ввод\_функции)  & \centering SGF1 & \centering Не предусмотрено \\ Предусмотрено & \centering -- & \centering \arraybackslash -- \\
\hline
\centering Блокировка от низкой изоляции (Блок\_от\_НИ)  & \centering SGF2 & \centering Не предусмотрено \\ Предусмотрено & \centering -- & \centering \arraybackslash -- \\
\hline
%===t1
\end{longtable}
\normalsize


%%%%%%%%%%%%%%%%%%%%%%%%%%%%%%%%%%%%%%%%%%%%%%%%%%%%%%%%%% ЛО ГЗсигн %%%%%%%%%%%%%%%%%%%%%%%%%%%%%%%%%%%%%%%%%%%%%%%%%%%%%%%%%%%%%%%%%%
\par

В таблице \ref{logzsign:tbl1} приведены параметры, необходимые для настройки ЛО ГЗсигн.

\footnotesize
\begin{longtable}{|>{\centering\arraybackslash}m{5.3cm}|>{\centering\arraybackslash}m{3.3cm}|>{\centering\arraybackslash}m{4.2cm}|>{\centering\arraybackslash}m{1.8cm}|>{\centering\arraybackslash}m{1cm}|}
\caption{Параметры для настройки логики отключения при срабатывании сигнального контакта газового реле трансформатора\hfill\vspace{-0.5\baselineskip}}\label{logzsign:tbl1}\\ 
\hline
\rowcolor{gray!20}
Параметр (Параметр на ИЧМ) & Условное обозначение на схеме & Значение/ Диапазон & Единица измерения & Шаг \\ 
\hline
\endfirsthead
\caption*{\hspace{3pt}\emph{Продолжение таблицы \ref{logzsign:tbl1}\hfill\vspace{-0.5\baselineskip}}} \\ % сделано по ГОСТ 2.105 п.6.8.7
\hline
\rowcolor{gray!20}
Параметр (Параметр на ИЧМ) & Условное обозначение на схеме & Значение/ Диапазон & Единица измерения & Шаг \\ 
%\hline
\endhead
%\hline
\endfoot
%\hline
\endlastfoot
%==+t1*PTRC1|TALMGASLGC
\multicolumn{5}{|c|}{ ЛО } \\ \hline 
\centering Ввод функции в работу (Ввод\_функции)  & \centering SGF1 & \centering Не предусмотрено \\ Предусмотрено & \centering -- & \centering \arraybackslash -- \\
\hline
\centering Выдержка времени срабатывания на блокировку (Тбл)  & \centering T1 & \centering 0...30000 & \centering мс & \centering \arraybackslash 10 \\
\hline
\centering Блокировка от низкой изоляции (Блок\_от\_НИ)  & \centering SGF2 & \centering Не предусмотрено \\ Предусмотрено & \centering -- & \centering \arraybackslash -- \\
\hline
%===t1
\end{longtable}
\normalsize

%%%%%%%%%%%%%%%%%%%%%%%%%%%%%%%%%%%%%%%%%%%%%%%%%%%%%%%%%% ЛО ГЗ РПН %%%%%%%%%%%%%%%%%%%%%%%%%%%%%%%%%%%%%%%%%%%%%%%%%%%%%%%%%%%%%%%%%%


\par

В таблице \ref{logzrpn:tbl1} приведены параметры, необходимые для настройки ЛО ГЗ РПН.

\footnotesize
\begin{longtable}{|>{\centering\arraybackslash}m{5.3cm}|>{\centering\arraybackslash}m{3.3cm}|>{\centering\arraybackslash}m{4.2cm}|>{\centering\arraybackslash}m{1.8cm}|>{\centering\arraybackslash}m{1cm}|}
\caption{Параметры для настройки логики отключения при срабатывании струйного реле РПН трансформатора\hfill\vspace{-0.5\baselineskip}}\label{logzrpn:tbl1}\\ 
\hline
\rowcolor{gray!20}
Параметр (Параметр на ИЧМ) & Условное обозначение на схеме & Значение/ Диапазон & Единица измерения & Шаг \\ 
\hline
\endfirsthead
\caption*{\hspace{3pt}\emph{Продолжение таблицы \ref{logzrpn:tbl1}\hfill\vspace{-0.5\baselineskip}}} \\ % сделано по ГОСТ 2.105 п.6.8.7
\hline
\rowcolor{gray!20}
Параметр (Параметр на ИЧМ) & Условное обозначение на схеме & Значение/ Диапазон & Единица измерения & Шаг \\ 
%\hline
\endhead
%\hline
\endfoot
%\hline
\endlastfoot
%==+t1*PTRC1|TLTCGASLGC
\multicolumn{5}{|c|}{ БУ } \\ \hline 
\centering Выдержка времени срабатывания на блокировку (Тбл)  & \centering T1 & \centering 0...30000 & \centering мс & \centering \arraybackslash 10 \\
\hline
\multicolumn{5}{|c|}{ ЛО } \\ \hline 
\centering Ввод функции в работу (Ввод\_функции)  & \centering SGF1 & \centering Не предусмотрено \\ Предусмотрено & \centering -- & \centering \arraybackslash -- \\
\hline
\centering Блокировка от низкой изоляции (Блок\_от\_НИ)  & \centering SGF2 & \centering Не предусмотрено \\ Предусмотрено & \centering -- & \centering \arraybackslash -- \\
\hline
%===t1
\end{longtable}
\normalsize



%%%%%%%%%%%%%%%%%%%%%%%%%%%%%%%%%%%%%%%%%%%%%%%%%%%%%%%%%% ЛО ТЗ %%%%%%%%%%%%%%%%%%%%%%%%%%%%%%%%%%%%%%%%%%%%%%%%%%%%%%%%%%%%%%%%%%
\par

В таблице \ref{lotz:tbl1} приведены параметры, необходимые для настройки ЛО ТЗ.
\footnotesize
\begin{longtable}{|>{\centering\arraybackslash}m{5.3cm}|>{\centering\arraybackslash}m{3.3cm}|>{\centering\arraybackslash}m{4.2cm}|>{\centering\arraybackslash}m{1.8cm}|>{\centering\arraybackslash}m{1cm}|}
\caption{Общие параметры для настройки ЛО ТЗ\hfill\vspace{-0.5\baselineskip}}\label{lotz:tbl1}\\ 
\hline
\rowcolor{gray!20}
Параметр (Параметр на ИЧМ) & Условное обозначение на схеме & Значение/ Диапазон & Единица измерения & Шаг \\ 
\hline
\endfirsthead
%\caption{\emph{Продолжение\hfill\vspace{-0.5\baselineskip}}} \\ % Отступ обозначения таблицы от таблицы и выравнивание надписи слева
\caption*{\hspace{3pt}\emph{Продолжение таблицы \ref{lotz:tbl1}\hfill\vspace{-0.5\baselineskip}}} \\ % сделано по ГОСТ 2.105 п.6.8.7
\hline
\rowcolor{gray!20}
Параметр (Параметр на ИЧМ) & Условное обозначение на схеме & Значение/ Диапазон & Единица измерения & Шаг \\ 
%\hline
\endhead
%\hline
\endfoot
%\hline
\endlastfoot
%==+t1*LLN0|APTTECHLGC
\multicolumn{5}{|c|}{ БУ } \\ \hline 
\centering Выдержка времени срабатывания на блокировку (Тбл)  & \centering T1 & \centering 0...30000 & \centering мс & \centering \arraybackslash 10 \\
\hline
\multicolumn{5}{|c|}{ ЛО ДТм } \\ \hline 
\centering Ввод функции в работу (Ввод\_функции)  & \centering SGF1 & \centering Не предусмотрено \\ Предусмотрено & \centering -- & \centering \arraybackslash -- \\
\hline
\centering Блокировка от низкой изоляции (Блок\_от\_НИ)  & \centering SGF2 & \centering Не предусмотрено \\ Предусмотрено & \centering -- & \centering \arraybackslash -- \\
\hline
\multicolumn{5}{|c|}{ ЛО ДТо } \\ \hline 
\centering Ввод функции в работу (Ввод\_функции)  & \centering SGF1 & \centering Не предусмотрено \\ Предусмотрено & \centering -- & \centering \arraybackslash -- \\
\hline
\centering Блокировка от низкой изоляции (Блок\_от\_НИ)  & \centering SGF2 & \centering Не предусмотрено \\ Предусмотрено & \centering -- & \centering \arraybackslash -- \\
\hline
\multicolumn{5}{|c|}{ ЛО РД } \\ \hline 
\centering Ввод функции в работу (Ввод\_функции)  & \centering SGF1 & \centering Не предусмотрено \\ Предусмотрено & \centering -- & \centering \arraybackslash -- \\
\hline
\centering Блокировка от низкой изоляции (Блок\_от\_НИ)  & \centering SGF2 & \centering Не предусмотрено \\ Предусмотрено & \centering -- & \centering \arraybackslash -- \\
\hline
%===t1
\end{longtable}
\normalsize


%%%%%%%%%%%%%%%%%%%%%%%%%%%%%%%%%%%%%%%%%%%%%%%%%%%%%%%%%% ЛО ТС %%%%%%%%%%%%%%%%%%%%%%%%%%%%%%%%%%%%%%%%%%%%%%%%%%%%%%%%%%%%%%%%%%
\par

В таблице \ref{lotspk:tbl1} приведены параметры, необходимые для настройки ЛО ТС.

\footnotesize
\begin{longtable}{|>{\centering\arraybackslash}m{5.3cm}|>{\centering\arraybackslash}m{3.3cm}|>{\centering\arraybackslash}m{4.2cm}|>{\centering\arraybackslash}m{1.8cm}|>{\centering\arraybackslash}m{1cm}|}
\caption{Параметры для настройки ЛО ПК\hfill\vspace{-0.5\baselineskip}}\label{lotspk:tbl1}\\ 
\hline
\rowcolor{gray!20}
Параметр (Параметр на ИЧМ) & Условное обозначение на схеме & Значение/ Диапазон & Единица измерения & Шаг \\ 
\hline
\endfirsthead
\caption*{\hspace{3pt}\emph{Продолжение таблицы \ref{lotspk:tbl1}\hfill\vspace{-0.5\baselineskip}}} \\ % сделано по ГОСТ 2.105 п.6.8.7
\hline
\rowcolor{gray!20}
Параметр (Параметр на ИЧМ) & Условное обозначение на схеме & Значение/ Диапазон & Единица измерения & Шаг \\ 
%\hline
\endhead
%\hline
\endfoot
%\hline
\endlastfoot
%==+t1*PRVLVPTRC1|ALMTECHLGC_UIRZ
\multicolumn{5}{|c|}{ ЛО ПК } \\ \hline 
\centering Ввод функции в работу (Ввод\_функции)  & \centering SGF1 & \centering Не предусмотрено \\ Предусмотрено & \centering -- & \centering \arraybackslash -- \\
\hline
\multicolumn{5}{|c|}{ ЛО ОК } \\ \hline 
\centering Ввод функции в работу (Ввод\_функции)  & \centering SGF1 & \centering Не предусмотрено \\ Предусмотрено & \centering -- & \centering \arraybackslash -- \\
\hline
\multicolumn{5}{|c|}{ ЛО ДУм расширителя } \\ \hline 
\centering Ввод функции в работу (Ввод\_функции)  & \centering SGF1 & \centering Не предусмотрено \\ Предусмотрено & \centering -- & \centering \arraybackslash -- \\
\hline
%===t1
\end{longtable}
\normalsize


%%%%%%%%%%%%%%%%%%%%%%%%%%%%%%%%%%%%%%%%%%%%%%%%%%%%%%%%%% ЛО Т %%%%%%%%%%%%%%%%%%%%%%%%%%%%%%%%%%%%%%%%%%%%%%%%%%%%%%%%%%%%%%%%%%
\par

В таблице \ref{loobsh:tbl1} приведены параметры, необходимые для настройки ЛО Т.

\footnotesize
\begin{longtable}{|>{\centering\arraybackslash}m{5.3cm}|>{\centering\arraybackslash}m{3.3cm}|>{\centering\arraybackslash}m{4.2cm}|>{\centering\arraybackslash}m{1.8cm}|>{\centering\arraybackslash}m{1cm}|}
\caption{Параметры для настройки функции <<ЛО>>\hfill\vspace{-0.5\baselineskip}}\label{loobsh:tbl1}\\ 
\hline
\rowcolor{gray!20}
Параметр (Параметр на ИЧМ) & Условное обозначение на схеме & Значение/ Диапазон & Единица измерения & Шаг \\ 
\hline
\endfirsthead
\caption*{\hspace{3pt}\emph{Продолжение таблицы \ref{loobsh:tbl1}\hfill\vspace{-0.5\baselineskip}}} \\ % сделано по ГОСТ 2.105 п.6.8.7
\hline
\rowcolor{gray!20}
Параметр (Параметр на ИЧМ) & Условное обозначение на схеме & Значение/ Диапазон & Единица измерения & Шаг \\ 
%\hline
\endhead
%\hline
\endfoot
%\hline
\endlastfoot
%==+t1*PTRC1|TOFFLVLGC
\multicolumn{5}{|c|}{ ЛО } \\ \hline 
\centering Ввод функции в работу (Ввод\_функции)  & \centering SGF1 & \centering Не предусмотрено \\ Предусмотрено & \centering -- & \centering \arraybackslash -- \\
\hline
\multicolumn{5}{|c|}{ ЗАПВ } \\ \hline 
\centering Ввод функции в работу (Ввод\_функции)  & \centering SGF1 & \centering Не предусмотрено \\ Предусмотрено & \centering -- & \centering \arraybackslash -- \\
\hline
\centering Режим запрета АПВ при срабатывании МТЗ 2 ступени (ЗАПВ\_МТЗ\_2ст)  & \centering SGF2 & \centering Не предусмотрено \\ Предусмотрено & \centering -- & \centering \arraybackslash -- \\
\hline
\centering Режим запрета АПВ при срабатывании МТЗ 3 ступени (ЗАПВ\_МТЗ\_3ст)  & \centering SGF3 & \centering Не предусмотрено \\ Предусмотрено & \centering -- & \centering \arraybackslash -- \\
\hline
\multicolumn{5}{|c|}{ ЗАВР } \\ \hline 
\centering Ввод функции в работу (Ввод\_функции)  & \centering SGF1 & \centering Не предусмотрено \\ Предусмотрено & \centering -- & \centering \arraybackslash -- \\
\hline
\centering Режим запрета АВР при срабатывании МТЗ 2 ступени (ЗАВР\_МТЗ\_2ст)  & \centering SGF2 & \centering Не предусмотрено \\ Предусмотрено & \centering -- & \centering \arraybackslash -- \\
\hline
\centering Режим запрета АВР при срабатывании МТЗ 3 ступени (ЗАВР\_МТЗ\_3ст)  & \centering SGF3 & \centering Не предусмотрено \\ Предусмотрено & \centering -- & \centering \arraybackslash -- \\
\hline
%===t1
\end{longtable}
\normalsize

%%%%%%%%%%%%%%%%%%%%%%%%%%%%%%%%%%%%%%%%%%%%%%%%%%%%%%%%%% ЛО ВН %%%%%%%%%%%%%%%%%%%%%%%%%%%%%%%%%%%%%%%%%%%%%%%%%%%%%%%%%%%%%%%%%%
\par
В таблице \ref{lovn:tbl1} приведены параметры, необходимые для настройки ЛО ВН.

\footnotesize
\begin{longtable}{|>{\centering\arraybackslash}m{5.3cm}|>{\centering\arraybackslash}m{3.3cm}|>{\centering\arraybackslash}m{4.2cm}|>{\centering\arraybackslash}m{1.8cm}|>{\centering\arraybackslash}m{1cm}|}
\caption{Параметры для настройки ЛО ВН\hfill\vspace{-0.5\baselineskip}}\label{lovn:tbl1}\\ 
\hline
\rowcolor{gray!20}
Параметр (Параметр на ИЧМ) & Условное обозначение на схеме & Значение/ Диапазон & Единица измерения & Шаг \\ 
\hline
\endfirsthead
\caption*{\hspace{3pt}\emph{Продолжение таблицы \ref{lovn:tbl1}\hfill\vspace{-0.5\baselineskip}}} \\ % сделано по ГОСТ 2.105 п.6.8.7
\hline
\rowcolor{gray!20}
Параметр (Параметр на ИЧМ) & Условное обозначение на схеме & Значение/ Диапазон & Единица измерения & Шаг \\ 
%\hline
\endhead
%\hline
\endfoot
%\hline
\endlastfoot
%==+t1*HVCBPTRC1|T_HVTCBOFF
\multicolumn{5}{|c|}{ ЛО } \\ \hline 
\centering Ввод функции в работу (Ввод\_функции)  & \centering SGF1 & \centering Не предусмотрено \\ Предусмотрено & \centering -- & \centering \arraybackslash -- \\
\hline
\centering Длительность импульса (Тимп)  & \centering T1 & \centering 50...60000 & \centering мс & \centering \arraybackslash 10 \\
\hline
%===t1
\end{longtable}
\normalsize


%%%%%%%%%%%%%%%%%%%%%%%%%%%%%%%%%%%%%%%%%%%%%%%%%%%%%%%%%% ЛО НН %%%%%%%%%%%%%%%%%%%%%%%%%%%%%%%%%%%%%%%%%%%%%%%%%%%%%%%%%%%%%%%%%%
\footnotesize
В таблице \ref{lonn:tbl_lov} приведены параметры, необходимые для настройки функции.

\small
\begin{longtable}{|>{\centering\arraybackslash}m{5cm}|>{\centering\arraybackslash}m{4cm}|>{\centering\arraybackslash}m{4cm}|>{\centering\arraybackslash}m{2cm}|>{\centering\arraybackslash}m{1cm}|}
\caption{Параметры для настройки логики отключения выключателя\hfill\vspace{-0.5\baselineskip}}\label{lonn:tbl_lov}\\ 
\hline
\rowcolor{gray!20}
Параметр (Параметр на ИЧМ) & Условное обозначение на схеме & Значение/ Диапазон & Единица измерения & Шаг \\ 
\hline
\endfirsthead
\caption*{\hspace{3pt}\emph{Продолжение таблицы \ref{lonn:tbl_lov}\hfill\vspace{-0.5\baselineskip}}} \\ % сделано по ГОСТ 2.105 п.6.8.7
\hline
\rowcolor{gray!20}
Параметр (Параметр на ИЧМ) & Условное обозначение на схеме & Значение/ Диапазон & Единица измерения & Шаг \\ 
%\hline
\endhead
%\hline
\endfoot
%\hline
\endlastfoot
%==+t1*LVCBPTRC1|LVTCBOFF
\multicolumn{5}{|c|}{ ЛО } \\ \hline 
\centering Ввод функции в работу (Ввод\_функции)  & \centering SGF1 & \centering Не предусмотрено \\ Предусмотрено & \centering -- & \centering \arraybackslash -- \\
\hline
\centering Длительность импульса (Тимп)  & \centering T1 & \centering 50...60000 & \centering мс & \centering \arraybackslash 10 \\
\hline
\multicolumn{5}{|c|}{ ЗАПВ } \\ \hline 
\centering Ввод функции в работу (Ввод\_функции)  & \centering SGF1 & \centering Не предусмотрено \\ Предусмотрено & \centering -- & \centering \arraybackslash -- \\
\hline
\multicolumn{5}{|c|}{ ЗАВР } \\ \hline 
\centering Ввод функции в работу (Ввод\_функции)  & \centering SGF1 & \centering Не предусмотрено \\ Предусмотрено & \centering -- & \centering \arraybackslash -- \\
\hline
%===t1
\end{longtable}
\normalsize

%%%%%%%%%%%%%%%%%%%%%%%%%%%%%%%%%%%%%%%%%%%%%%%%%%%%%%%%%% УРОВ %%%%%%%%%%%%%%%%%%%%%%%%%%%%%%%%%%%%%%%%%%%%%%%%%%%%%%%%%%%%%%%%%%
\par
В таблице \ref{urov35:tbl1} приведены параметры, необходимые для настройки УРОВ.

\footnotesize
\begin{longtable}{|>{\centering\arraybackslash}m{5.3cm}|>{\centering\arraybackslash}m{3.3cm}|>{\centering\arraybackslash}m{4.2cm}|>{\centering\arraybackslash}m{1.8cm}|>{\centering\arraybackslash}m{1cm}|}
\caption{Параметры для настройки УРОВ\hfill\vspace{-0.5\baselineskip}}\label{urov35:tbl1}\\ 
\hline
\rowcolor{gray!20}
Параметр (Параметр на ИЧМ) & Условное обозначение на схеме & Значение/ Диапазон & Единица измерения & Шаг \\ 
\hline
\endfirsthead
\caption*{\hspace{3pt}\emph{Продолжение таблицы \ref{urov35:tbl1}\hfill\vspace{-0.5\baselineskip}}} \\ % сделано по ГОСТ 2.105 п.6.8.7
\hline
\rowcolor{gray!20}
Параметр (Параметр на ИЧМ) & Условное обозначение на схеме & Значение/ Диапазон & Единица измерения & Шаг \\ 
%\hline
\endhead
%\hline
\endfoot
%\hline
\endlastfoot
%==+t1*GENRBRF1|T_TPBRF
\multicolumn{5}{|c|}{ УРОВ } \\ \hline 
\centering Ввод функции в работу (Ввод\_функции)  & \centering SGF1 & \centering Не предусмотрено \\ Предусмотрено & \centering -- & \centering \arraybackslash -- \\
\hline
\centering Ток срабатывания (Iср)  & \centering I> & \centering 0,05...0,50 & \centering о.е. & \centering \arraybackslash 0,01 \\
\hline
\centering Ускорение при блокировке отключения В (Уск\_при\_блк\_откл\_В)  & \centering SGF2 & \centering Не предусмотрено \\ Предусмотрено & \centering -- & \centering \arraybackslash -- \\
\hline
\centering УРОВ с подхватом по току (Подхват\_по\_току)  & \centering SGF3 & \centering Не предусмотрено \\ Предусмотрено & \centering -- & \centering \arraybackslash -- \\
\hline
\centering Контроль по току при действии "на себя" (Контр\_тока\_на\_себя)  & \centering SGF6 & \centering Не предусмотрено \\ Предусмотрено по внутр. ПО \\ Предусмотрено по внеш. ПО & \centering -- & \centering \arraybackslash -- \\
\hline
\centering Контроль ЭМО (Контроль\_ЭМО)  & \centering SGF4 & \centering Не предусмотрено \\ Предусмотрено по ЭМО1 \\ Предусмотрено по ЭМО1 и ЭМО2 & \centering -- & \centering \arraybackslash -- \\
\hline
\centering Действие внешнего УРОВ на вышестоящий выключатель (Действ\_на\_выш\_выкл)  & \centering SGF5 & \centering Не предусмотрено \\ Предусмотрено & \centering -- & \centering \arraybackslash -- \\
\hline
\centering Выдержка времени срабатывания (Tср)  & \centering T1 & \centering 0...1000 & \centering мс & \centering \arraybackslash 10 \\
\hline
%===t1
\end{longtable}
\normalsize



%%%%%%%%%%%%%%%%%%%%%%%%%%%%%%%%%%%%%%%%%%%%%%%%%%%%%%%%%% КРВ %%%%%%%%%%%%%%%%%%%%%%%%%%%%%%%%%%%%%%%%%%%%%%%%%%%%%%%%%%%%%%%%%%
\par
В таблице \ref{krv:tbl1} приведены параметры, необходимые для настройки КРВ.

\footnotesize
\begin{longtable}{|>{\centering\arraybackslash}m{5.3cm}|>{\centering\arraybackslash}m{3.3cm}|>{\centering\arraybackslash}m{4.2cm}|>{\centering\arraybackslash}m{1.8cm}|>{\centering\arraybackslash}m{1cm}|}
\caption{Параметры для настройки КРВ\hfill\vspace{-0.5\baselineskip}}\label{krv:tbl1}\\ 
\hline
\rowcolor{gray!20} Параметр (Параметр на ИЧМ) & Условное обозначение на схеме & Значение/ Диапазон & Единица измерения & Шаг \\ 
\hline
\endfirsthead
\caption*{\hspace{3pt}\emph{Продолжение таблицы \ref{krv:tbl1}\hfill\vspace{-0.5\baselineskip}}} \\ % сделано по ГОСТ 2.105 п.6.8.7
\hline
\rowcolor{gray!20} Параметр (Параметр на ИЧМ) & Условное обозначение на схеме & Значение/ Диапазон & Единица измерения & Шаг \\ 
%\hline
\endhead
%\hline
\endfoot
%\hline
\endlastfoot 
%==+t1*CLS|CLS
\centering Ввод функции в работу (Ввод\_функции)  & \centering SGF1 & \centering Не предусмотрено \\ Предусмотрено & \centering -- & \centering \arraybackslash -- \\
\hline
\centering Номинальный ток выключателя (паспорт) (Iном\_В\_пасп)  & \centering -- & \centering 0...5000 & \centering А & \centering \arraybackslash 1 \\
\hline
\centering Номинальный ток отключения выключателя (паспорт) (Iном\_откл\_В\_пасп)  & \centering -- & \centering 0...40000 & \centering А & \centering \arraybackslash 1 \\
\hline
\centering Допустимое количество циклов при номинальном токе выключателя (паспорт) (КРВпасп\_Iном)  & \centering -- & \centering 0...200000 & \centering -- & \centering \arraybackslash 1 \\
\hline
\centering Допустимое количество циклов при номинальном токе отключения выключателя (паспорт) (КРВпасп\_Iном\_откл)  & \centering -- & \centering 0...500 & \centering -- & \centering \arraybackslash 1 \\
\hline
\centering Допустимый механический ресурс выключателя (паспорт) (МРВпасп)  & \centering -- & \centering 0...200000 & \centering -- & \centering \arraybackslash 1 \\
\hline
\centering Начальное значение коммутационного ресурса выключателя (Нач\_знач\_КРВ)  & \centering -- & \centering 0,0...100,0 & \centering \% & \centering \arraybackslash 0,1 \\
\hline
\centering Величина срабатывания коммутационного ресурса выключателя (КРВсраб)  & \centering -- & \centering 0,0...100,0 & \centering \% & \centering \arraybackslash 0,1 \\
\hline
\centering Начальное значение механического ресурса выключателя (Нач\_знач\_МРВ)  & \centering -- & \centering 0...200000 & \centering -- & \centering \arraybackslash 1 \\
\hline
\centering Контроль КРВ по превышению МРВ (Контроль\_КРВ\_от\_МРВ)  & \centering SGF2 & \centering Не предусмотрено \\ Предусмотрено & \centering -- & \centering \arraybackslash -- \\
\hline
\centering Полное время отключения выключателя (Tмакс\_В)  & \centering -- & \centering 10...3000 & \centering мс & \centering \arraybackslash 10 \\
\hline
%===t1

\end{longtable}
\normalsize

%%%%%%%%%%%%%%%%%%%%%%%%%%%%%%%%%%%%%%%%%%%%%%%%%%%%%%%%%% КСВ %%%%%%%%%%%%%%%%%%%%%%%%%%%%%%%%%%%%%%%%%%%%%%%%%%%%%%%%%%%%%%%%%%
\par
В таблице \ref{ksv:tbl1} приведены параметры, необходимые для настройки КСВ.

\footnotesize
\begin{longtable}{|>{\centering\arraybackslash}m{5.3cm}|>{\centering\arraybackslash}m{3.3cm}|>{\centering\arraybackslash}m{4.2cm}|>{\centering\arraybackslash}m{1.8cm}|>{\centering\arraybackslash}m{1cm}|}
\caption{Параметры для настройки КСВ\hfill\vspace{-0.5\baselineskip}}\label{ksv:tbl1}\\ 
\hline
\rowcolor{gray!20}
Параметр (Параметр на ИЧМ) & Условное обозначение на схеме & Значение/ Диапазон & Единица измерения & Шаг \\ 
\hline
\endfirsthead
\caption*{\hspace{3pt}\emph{Продолжение таблицы \ref{ksv:tbl1}\hfill\vspace{-0.5\baselineskip}}} \\ % сделано по ГОСТ 2.105 п.6.8.7
\hline
\rowcolor{gray!20}
Параметр (Параметр на ИЧМ) & Условное обозначение на схеме & Значение/ Диапазон & Единица измерения & Шаг \\ 
%\hline
\endhead
%\hline
\endfoot
%\hline
\endlastfoot
%==+t1*RCBF1|T_LVCBSUP
\centering Ввод функции в работу (Ввод\_функции)  & \centering SGF1 & \centering Не предусмотрено \\ Предусмотрено & \centering -- & \centering \arraybackslash -- \\
\hline
\centering Контроль ОТ цепей ЭМВ, ЭМО1 и ЭМО2 (Контроль\_ОТ\_ЭМ)  & \centering SGF5 & \centering Не предусмотрено \\ ЭМВ и ЭМО1 \\ ЭМВ, ЭМО1 и ЭМО2 & \centering -- & \centering \arraybackslash -- \\
\hline
\centering Блокировка управления при снижении уровня изоляции В (Блок\_упр\_КИ\_В)  & \centering SGF8 & \centering Не предусмотрено \\ От аварийного \\ От аварийного и низкого & \centering -- & \centering \arraybackslash -- \\
\hline
\centering Блокировка включения при низком уровне изоляции В (Блок\_вкл\_низ\_изол\_В)  & \centering SGF2 & \centering Не предусмотрено \\ Предусмотрено & \centering -- & \centering \arraybackslash -- \\
\hline
\centering Блокировка включения при неисправности положения В (Блок\_вкл\_полож\_В)  & \centering SGF3 & \centering Не предусмотрено \\ Предусмотрено & \centering -- & \centering \arraybackslash -- \\
\hline
\centering Блокировка включения при превышения ресурса В (Блок\_вкл\_ресурса\_В)  & \centering SGF4 & \centering Не предусмотрено \\ Предусмотрено & \centering -- & \centering \arraybackslash -- \\
\hline
\centering Контроль ЭМВ, ЭМО1 и ЭМО2 при формировании неисправности цепей ЭМУ (Контроль\_ЭМ)  & \centering SGF6 & \centering Не предусмотрено \\ ЭМВ и ЭМО1 \\ ЭМВ, ЭМО1 и ЭМО2 & \centering -- & \centering \arraybackslash -- \\
\hline
\centering Разрешение сброса "РФК" от кнопки (Контроль\_кнопки)  & \centering SGF7 & \centering Не предусмотрено \\ Предусмотрено & \centering -- & \centering \arraybackslash -- \\
\hline
\centering Выдержка времени ожидания готовности привода выключателя (Tгот\_в)  & \centering T1 & \centering 0...100000 & \centering мс & \centering \arraybackslash 10 \\
\hline
\centering Выдержка времени срабатывания неисправности цепей ЭМУ (Tдлит\_неисп\_ЭМУ)  & \centering T2 & \centering 0...10000 & \centering мс & \centering \arraybackslash 10 \\
\hline
\centering Выдержка времени срабатывания защит электромагнитов (Tдлит\_ЭМ)  & \centering T3 & \centering 0...10000 & \centering мс & \centering \arraybackslash 10 \\
\hline
%===t1
\end{longtable}
\normalsize


%%%%%%%%%%%%%%%%%%%%%%%%%%%%%%%%%%%%%%%%%%%%%%%%%%%%%%%%%% КП %%%%%%%%%%%%%%%%%%%%%%%%%%%%%%%%%%%%%%%%%%%%%%%%%%%%%%%%%%%%%%%%%%
\par
В таблице \ref{kp:tbl_gen} приведены параметры, необходимые для настройки КП. 

\footnotesize
\begin{longtable}{|>{\centering\arraybackslash}m{5.3cm}|>{\centering\arraybackslash}m{3.3cm}|>{\centering\arraybackslash}m{4.2cm}|>{\centering\arraybackslash}m{1.8cm}|>{\centering\arraybackslash}m{1cm}|}
\caption{Параметры для настройки КП\hfill\vspace{-0.5\baselineskip}}\label{kp:tbl_gen}\\ 
\hline
\rowcolor{gray!20}
Параметр (Параметр на ИЧМ) & Условное обозначение на схеме & Значение/ Диапазон & Единица измерения & Шаг \\ 
\hline
\endfirsthead
\caption*{\hspace{3pt}\emph{Продолжение таблицы \ref{kp:tbl_gen}\hfill\vspace{-0.5\baselineskip}}} \\ % сделано по ГОСТ 2.105 п.6.8.7
\hline
\rowcolor{gray!20}
Параметр (Параметр на ИЧМ) & Условное обозначение на схеме & Значение/ Диапазон & Единица измерения & Шаг \\ 
%\hline
\endhead
%\hline
\endfoot
%\hline
\endlastfoot
%==+t1*LLN0|T_SWCTRL
\multicolumn{5}{|c|}{ УВ } \\ \hline 
\centering Ввод функции в работу (Ввод\_функции)  & \centering SGF1 & \centering Не предусмотрено \\ Предусмотрено & \centering -- & \centering \arraybackslash -- \\
\hline
\centering Блокировка включения от аварийного отключения В (Бл\_вкл\_от\_авар\_откл)  & \centering SGF2 & \centering Не предусмотрено \\ Предусмотрено & \centering -- & \centering \arraybackslash -- \\
\hline
\centering Допустимое время переключения выключателя (Tперекл)  & \centering T1 & \centering 10...10000 & \centering мс & \centering \arraybackslash 10 \\
\hline
\centering Выдержка времени подавления выдачи положения "Не определено" (Tблк)  & \centering T2 & \centering 10...10000 & \centering мс & \centering \arraybackslash 10 \\
\hline
\multicolumn{5}{|c|}{ КП } \\ \hline 
\centering Ввод функции в работу (Ввод\_функции)  & \centering SGF1 & \centering Не предусмотрено \\ Предусмотрено & \centering -- & \centering \arraybackslash -- \\
\hline
%===t1
\end{longtable}
\normalsize


%%%%%%%%%%%%%%%%%%%%%%%%%%%%%%%%%%%%%%%%%%%%%%%%%%%%%%%%%% УВ %%%%%%%%%%%%%%%%%%%%%%%%%%%%%%%%%%%%%%%%%%%%%%%%%%%%%%%%%%%%%%%%%%
\par
В таблице \ref{uv:tbl_gen} приведены параметры, необходимые для настройки УВ. 

\footnotesize
\begin{longtable}{|>{\centering\arraybackslash}m{5.3cm}|>{\centering\arraybackslash}m{3.3cm}|>{\centering\arraybackslash}m{4.2cm}|>{\centering\arraybackslash}m{1.8cm}|>{\centering\arraybackslash}m{1cm}|}
\caption{Параметры для настройки УВ\hfill\vspace{-0.5\baselineskip}}\label{uv:tbl_gen}\\ 
\hline
\rowcolor{gray!20}
Параметр (Параметр на ИЧМ) & Условное обозначение на схеме & Значение/ Диапазон & Единица измерения & Шаг \\ 
\hline
\endfirsthead
\caption*{\hspace{3pt}\emph{Продолжение таблицы \ref{uv:tbl_gen}\hfill\vspace{-0.5\baselineskip}}} \\ % сделано по ГОСТ 2.105 п.6.8.7
\hline
\rowcolor{gray!20}
Параметр (Параметр на ИЧМ) & Условное обозначение на схеме & Значение/ Диапазон & Единица измерения & Шаг \\ 
%\hline
\endhead
%\hline
\endfoot
%\hline
\endlastfoot
%==+t1*CBCSWI1|T_HVBCTRL
\centering Ввод функции в работу (Ввод\_функции)  & \centering SGF1 & \centering Не предусмотрено \\ Предусмотрено & \centering -- & \centering \arraybackslash -- \\
\hline
%===t1
\end{longtable}
\normalsize


%%%%%%%%%%%%%%%%%%%%%%%%%%%%%%%%%%%%%%%%%%%%%%%%%%%%%%%%%% КА %%%%%%%%%%%%%%%%%%%%%%%%%%%%%%%%%%%%%%%%%%%%%%%%%%%%%%%%%%%%%%%%%%
\par
В таблице \ref{ka:tbl_gen} приведены параметры, необходимые для настройки КА.

\footnotesize
\begin{longtable}{|>{\centering\arraybackslash}m{5.3cm}|>{\centering\arraybackslash}m{3.3cm}|>{\centering\arraybackslash}m{4.2cm}|>{\centering\arraybackslash}m{1.8cm}|>{\centering\arraybackslash}m{1cm}|}
\caption{Параметры для настройки функционального блока «КА»\hfill\vspace{-0.5\baselineskip}}\label{ka:tbl_gen}\\ 
\hline
\rowcolor{gray!20}
Параметр (Параметр на ИЧМ) & Условное обозначение на схеме & Значение/ Диапазон & Единица измерения & Шаг \\ 
\hline
\endfirsthead
\caption*{\hspace{3pt}\emph{Продолжение таблицы \ref{ka:tbl_gen}\hfill\vspace{-0.5\baselineskip}}} \\ % сделано по ГОСТ 2.105 п.6.8.7
\hline
\rowcolor{gray!20}
Параметр (Параметр на ИЧМ) & Условное обозначение на схеме & Значение/ Диапазон & Единица измерения & Шаг \\ 
%\hline
\endhead
%\hline
\endfoot
%\hline
\endlastfoot
%==+t1*LLN0|T_SwitchDevice
\multicolumn{5}{|c|}{ В } \\ \hline 
\centering Ввод функции в работу (Ввод\_функции)  & \centering SGF1 & \centering Не предусмотрено \\ Предусмотрено & \centering -- & \centering \arraybackslash -- \\
\hline
\centering Режим отключения В (Режим\_откл)  & \centering SGF2 & \centering Длительный \\ Импульсный & \centering -- & \centering \arraybackslash -- \\
\hline
\centering Контроль работы ЭМО (Контр\_раб\_ЭМО)  & \centering SGF3 & \centering Не предусмотрено \\ Предусмотрено & \centering -- & \centering \arraybackslash -- \\
\hline
\centering Режим включения В (Режим\_вкл)  & \centering SGF4 & \centering Длительный \\ Импульсный & \centering -- & \centering \arraybackslash -- \\
\hline
\centering Контроль работы ЭМВ (Контр\_раб\_ЭМВ)  & \centering SGF5 & \centering Не предусмотрено \\ Предусмотрено & \centering -- & \centering \arraybackslash -- \\
\hline
\centering Длительность формирования сигнала о неисправном положении В (Tнеиспр\_В)  & \centering T1 & \centering 10...10000 & \centering мс & \centering \arraybackslash 10 \\
\hline
\centering Длительность формирования команды "Отключить В (реле)" (Тимп\_откл)  & \centering T2 & \centering 50...10000 & \centering мс & \centering \arraybackslash 10 \\
\hline
\centering Длительность формирования команды "Включить В (реле)" (Тимп\_вкл)  & \centering T3 & \centering 50...10000 & \centering мс & \centering \arraybackslash 10 \\
\hline
\centering Время для исключения "опрокидывания" В (Тпрод\_вкл)  & \centering T4 & \centering 0...1000 & \centering мс & \centering \arraybackslash 10 \\
\hline
\multicolumn{5}{|c|}{ КА } \\ \hline 
\centering Ввод функции в работу (Ввод\_функции)  & \centering SGF1 & \centering Не предусмотрено \\ Предусмотрено & \centering -- & \centering \arraybackslash -- \\
\hline
%===t1
\end{longtable}
\normalsize


%%%%%%%%%%%%%%%%%%%%%%%%%%%%%%%%%%%%%%%%%%%%%%%%%%%%%%%%%% ПС %%%%%%%%%%%%%%%%%%%%%%%%%%%%%%%%%%%%%%%%%%%%%%%%%%%%%%%%%%%%%%%%%%
\par
В таблице \ref{ps:tbl1} приведены параметры, необходимые для настройки ПС. 

\footnotesize
\begin{longtable}{|>{\centering\arraybackslash}m{5.3cm}|>{\centering\arraybackslash}m{3.3cm}|>{\centering\arraybackslash}m{4.2cm}|>{\centering\arraybackslash}m{1.8cm}|>{\centering\arraybackslash}m{1cm}|}
\caption{Параметры для настройки предупредительной сигнализации\hfill\vspace{-0.5\baselineskip}}\label{ps:tbl1}\\ 
\hline
\rowcolor{gray!20}
Параметр (Параметр на ИЧМ) & Условное обозначение на схеме & Значение/ Диапазон & Единица измерения & Шаг \\ 
\hline
\endfirsthead
\caption*{\hspace{3pt}\emph{Продолжение таблицы \ref{ps:tbl1}\hfill\vspace{-0.5\baselineskip}}} \\ % сделано по ГОСТ 2.105 п.6.8.7
\hline
\rowcolor{gray!20}
Параметр (Параметр на ИЧМ) & Условное обозначение на схеме & Значение/ Диапазон & Единица измерения & Шаг \\ 
%\hline
\endhead
%\hline
\endfoot
%\hline
\endlastfoot
%==+t1*CALH1|T_LVALH
\centering Контроль сигнала «ГЗ сигн» (Контр\_ГЗ\_сигн)  & \centering SGF1 & \centering Не предусмотрено \\ Предусмотрено & \centering -- & \centering \arraybackslash -- \\
\hline
\centering Контроль сигнала «Низ. изол. ГЗ» (Контр\_Низ\_изол\_ГЗ)  & \centering SGF2 & \centering Не предусмотрено \\ Предусмотрено & \centering -- & \centering \arraybackslash -- \\
\hline
\centering Контроль сигнала «ГЗ заблокирована» (Контр\_ГЗ\_блок)  & \centering SGF3 & \centering Не предусмотрено \\ Предусмотрено & \centering -- & \centering \arraybackslash -- \\
\hline
\centering Контроль сигнала «ТЗ сигн» (Контр\_ТЗ\_сигн)  & \centering SGF4 & \centering Не предусмотрено \\ Предусмотрено & \centering -- & \centering \arraybackslash -- \\
\hline
\centering Контроль сигнала «Низ. изол. ТЗ» (Контр\_Низ\_изол\_ТЗ)  & \centering SGF5 & \centering Не предусмотрено \\ Предусмотрено & \centering -- & \centering \arraybackslash -- \\
\hline
\centering Контроль сигнала «ТЗ заблокирована» (Контр\_ТЗ\_блок)  & \centering SGF6 & \centering Не предусмотрено \\ Предусмотрено & \centering -- & \centering \arraybackslash -- \\
\hline
\centering Контроль сигнала «ТС сигн» (Контр\_ТС\_сигн)  & \centering SGF7 & \centering Не предусмотрено \\ Предусмотрено & \centering -- & \centering \arraybackslash -- \\
\hline
\centering Контроль сигнала «ОТ сигн» (Контр\_ОТ\_сигн)  & \centering SGF8 & \centering Не предусмотрено \\ Предусмотрено & \centering -- & \centering \arraybackslash -- \\
\hline
\centering Контроль сигнала «ОТ НН сигн» (Контр\_ОТ\_НН\_сигн)  & \centering SGF9 & \centering Не предусмотрено \\ Предусмотрено & \centering -- & \centering \arraybackslash -- \\
\hline
\centering Контроль сигнала «Внеш. откл» (Контр\_Внеш\_откл)  & \centering SGF10 & \centering Не предусмотрено \\ Предусмотрено & \centering -- & \centering \arraybackslash -- \\
\hline
\centering Контроль сигнала «Вых. цепи разобраны» (Контр\_Вых\_цеп\_разоб)  & \centering SGF11 & \centering Не предусмотрено \\ Предусмотрено & \centering -- & \centering \arraybackslash -- \\
\hline
\centering Контроль сигнала «БИ выведены» (Контр\_БИ\_выведены)  & \centering SGF12 & \centering Не предусмотрено \\ Предусмотрено & \centering -- & \centering \arraybackslash -- \\
\hline
\centering Контроль сигнала «Прев. вр. пер. КА» (Контр\_Вр\_пркл\_КА)  & \centering SGF13 & \centering Не предусмотрено \\ Предусмотрено & \centering -- & \centering \arraybackslash -- \\
\hline
\centering Контроль общего внешнего сигнала (Контр\_общ\_внеш\_сигн)  & \centering SGF14 & \centering Не предусмотрено \\ Предусмотрено & \centering -- & \centering \arraybackslash -- \\
\hline
%===t1
\end{longtable}
\normalsize


%%%%%%%%%%%%%%%%%%%%%%%%%%%%%%%%%%%%%%%%%%%%%%%%%%%%%%%%%% СС %%%%%%%%%%%%%%%%%%%%%%%%%%%%%%%%%%%%%%%%%%%%%%%%%%%%%%%%%%%%%%%%%%
\par
В таблице \ref{ss:tbl1} приведены параметры, необходимые для настройки СС. 

\footnotesize
\begin{longtable}{|>{\centering\arraybackslash}m{5.3cm}|>{\centering\arraybackslash}m{3.3cm}|>{\centering\arraybackslash}m{4.2cm}|>{\centering\arraybackslash}m{1.8cm}|>{\centering\arraybackslash}m{1cm}|}
\caption{Параметры для настройки функции сборки сигналов\hfill\vspace{-0.5\baselineskip}}\label{ss:tbl1}\\ 
\hline
\rowcolor{gray!20}
Параметр (Параметр на ИЧМ) & Условное обозначение на схеме & Значение/ Диапазон & Единица измерения & Шаг \\ 
\hline
\endfirsthead
\caption*{\hspace{3pt}\emph{Продолжение таблицы \ref{ss:tbl1}\hfill\vspace{-0.5\baselineskip}}} \\ % сделано по ГОСТ 2.105 п.6.8.7
\hline
\rowcolor{gray!20}
Параметр (Параметр на ИЧМ) & Условное обозначение на схеме & Значение/ Диапазон & Единица измерения & Шаг \\ 
%\hline
\endhead
%\hline
\endfoot
%\hline
\endlastfoot
%==+t1*GAPC1|T_SignAssembly
\multicolumn{5}{|c|}{ CC } \\ \hline 
\centering Контроль положения SA1 (Контр\_полож\_SA1)  & \centering SGF1 & \centering Не предусмотрено \\ Предусмотрено & \centering -- & \centering \arraybackslash -- \\
\hline
\centering Контроль положения SA2 (Контр\_полож\_SA2)  & \centering SGF2 & \centering Не предусмотрено \\ Предусмотрено & \centering -- & \centering \arraybackslash -- \\
\hline
\centering Контроль положения SA3 (Контр\_полож\_SA3)  & \centering SGF3 & \centering Не предусмотрено \\ Предусмотрено & \centering -- & \centering \arraybackslash -- \\
\hline
\centering Контроль положения SA4 (Контр\_полож\_SA4)  & \centering SGF4 & \centering Не предусмотрено \\ Предусмотрено & \centering -- & \centering \arraybackslash -- \\
\hline
\centering Контроль положения SA5 (Контр\_полож\_SA5)  & \centering SGF5 & \centering Не предусмотрено \\ Предусмотрено & \centering -- & \centering \arraybackslash -- \\
\hline
\centering Контроль положения SG1 (Контр\_полож\_SG1)  & \centering SGF6 & \centering Не предусмотрено \\ Предусмотрено & \centering -- & \centering \arraybackslash -- \\
\hline
\centering Контроль положения SG2 (Контр\_полож\_SG2)  & \centering SGF7 & \centering Не предусмотрено \\ Предусмотрено & \centering -- & \centering \arraybackslash -- \\
\hline
\centering Контроль опер. тока цепей ГЗ (Контр\_ОТ\_ГЗ)  & \centering SGF8 & \centering Не предусмотрено \\ Предусмотрено & \centering -- & \centering \arraybackslash -- \\
\hline
\centering Контроль опер. тока цепей ТЗ, ТС (Контр\_ОТ\_ТЗ\_ТС)  & \centering SGF9 & \centering Не предусмотрено \\ Предусмотрено & \centering -- & \centering \arraybackslash -- \\
\hline
\centering Контроль опер. тока цепей сигнализации выключателя (Контр\_ОТ\_сигн\_В)  & \centering SGF10 & \centering Не предусмотрено \\ Предусмотрено & \centering -- & \centering \arraybackslash -- \\
\hline
\centering Контроль опер. тока цепей ЗДЗ НН (Контр\_ОТ\_ЗДЗ)  & \centering SGF11 & \centering Не предусмотрено \\ Предусмотрено & \centering -- & \centering \arraybackslash -- \\
\hline
\centering Контроль опер. тока цепей УРОВ НН (Контр\_ОТ\_УРОВ)  & \centering SGF12 & \centering Не предусмотрено \\ Предусмотрено & \centering -- & \centering \arraybackslash -- \\
\hline
\centering Контроль опер. тока цепей ОТ ИЭУ ТН (Контр\_ОТ\_ТН)  & \centering SGF13 & \centering Не предусмотрено \\ Предусмотрено & \centering -- & \centering \arraybackslash -- \\
\hline
%===t1
\end{longtable}
\normalsize